\section{Conclusion, Limitations, and Future Work}
\label{sec:Conclusion, Limiatations, and Future Work}

We introduced a scalable framework for recovering naturally occurring scientific-figure revisions from arXiv, GitHub, and TeX SE, instantiated as DaEdiTikZ, a large-scale real-world TikZ editing dataset. We further introduced the human-refined DaEdiTikZ-Bench and EdiTikZ, a family of 4--9B models trained with multi-task SFT and multi-reward RL. EdiTikZ-9B-RL leads automatic evaluation and reaches comparable human ratings to the strongest proprietary system. Post-training gains transfer to substantially more complex SPIQA and CharXiv figures, even beyond the 2K-token training horizon. Overall, naturally occurring revision trajectories provide effective supervision for small, open scientific-figure editing models competitive with much larger proprietary systems.

DaEdiTikZ inherits noise from automatically inferred instructions, including omissions and misinterpretations despite filtering and code grounding. Performance also degrades for long OOD generations, motivating post-training on more complex figures in the future. Evaluation in this regime is itself limited by synthetic instructions and potentially less reliable reference-free judging. Beyond these limitations, our visualization-language-agnostic revision-mining framework could extend to Matplotlib or LaTeX tables, while access to source programs could enable localized editing without full reconstruction. Revision trajectories could further support comparative VQA, retrieval, and representation learning, while helping to unify generation and editing within general-purpose scientific visualization models.